% =================================================
% Домашняя работа. Урок 1
% =================================================

\homeworktitle{Домашняя работа}{Цифры, числа и разряды}
\studentline

\begin{routebox}[Как выполнять]
  Сначала выполните задания 1--5. В задании 6 найдите и объясните ошибку.
  Задание 7 отмечено звёздочкой: оно требует перевести названные разрядные
  единицы в обычную запись числа.
\end{routebox}

\begin{homeworktask}[Цифры в записи числа]{1}
  Для каждого числа укажите количество цифр и перечислите цифры по порядку:
  \[
    36,\qquad 507,\qquad 70\,090.
  \]
  \answerspace{30mm}
\end{homeworktask}

\begin{homeworktask}[Назовите цифру разряда]{2}
  \begin{enumerate}
    \item цифру сотен числа \(48\,372\);
    \item цифру десятков тысяч числа \(605\,814\);
    \item цифру десятков числа \(90\,407\).
  \end{enumerate}
  \answerspace{26mm}
\end{homeworktask}

\begin{homeworktask}[Каково значение цифры?]{3}
  \begin{enumerate}
    \item цифры \(7\) в числе \(57\,214\);
    \item цифры \(3\) в числе \(203\,681\);
    \item цифры \(9\) в числе \(45\,092\).
  \end{enumerate}
  \answerspace{27mm}
\end{homeworktask}

\begin{homeworktask}[Разложите по разрядам]{4}
  \begin{enumerate}
    \item \(64\,205\);
    \item \(8\,090\);
    \item \(300\,407\).
  \end{enumerate}
  \answerspace{34mm}
\end{homeworktask}

\begin{homeworktask}[Соберите число]{5}
  \begin{enumerate}
    \item \(50\,000+6\,000+300+2\);
    \item \(90\,000+700+40\);
    \item \(400\,000+20\,000+5\,000+9\).
  \end{enumerate}
  \answerspace{30mm}
\end{homeworktask}

\begin{homeworktask}[Найдите ошибку]{6}
  Ученик написал:
  \[
    71\,304=70\,000+1\,000+3\,000+4.
  \]
  Объясните ошибку и запишите правильное разложение.
  \answerspace{36mm}
\end{homeworktask}

\begin{homeworktask}[Задание со звёздочкой]{7*}
  Какое число получится, если взять \(6\) десятков тысяч, \(34\) сотни
  и \(8\) единиц?

  \begin{hintbox}
    Переведите каждую часть в число:
    \(6\) десятков тысяч --- это \(6\cdot10\,000\),
    а \(34\) сотни --- это \(34\cdot100\).
  \end{hintbox}
  \answerspace{42mm}
\end{homeworktask}

\begin{selfcheck}
  \begin{itemize}[label=\(\square\),leftmargin=1.7em,itemsep=0.25em,topsep=0.2em]
    \item Я считаю все цифры, включая нули внутри записи.
    \item Я определяю разряды справа налево.
    \item В разложении каждая цифра имеет правильное значение.
    \item При сборке числа я сохраняю пустые разряды с помощью нулей.
  \end{itemize}
\end{selfcheck}
